% Vietnamese translation for OI Public Library
% Source-PDF: IOI中国国家候选队论文/2009/曹钦翔/从“k倍动态减法游戏”出发探究一类组合游戏问题.pdf
\documentclass[11pt]{article}
\usepackage{oipl-vn}

\newcommand{\Oh}{\mathcal{O}}
\newcommand{\NP}{\operatorname{NP}}
\newcommand{\SG}{\operatorname{SG}}
\newcommand{\mex}{\operatorname{mex}}
\newcommand{\xor}{\oplus}
\newcommand{\nimprod}{\otimes}
\newcommand{\lowbit}{\operatorname{lowbit}}
\newcommand{\lowbitF}{\operatorname{lowbit}_{F}}

\title{Từ trò chơi trừ động theo bội \(k\) đến một lớp bài toán trò chơi tổ hợp}
\author{Cao Qinxiang\\Trường Trung học số 2 trực thuộc Đại học Sư phạm Hoa Đông}
\date{2009}

\begin{document}
\maketitle

\origtitle{Từ trò chơi trừ động theo bội \(k\) đến một lớp bài toán trò chơi tổ hợp}
\sourcepath{IOI China candidate team papers / 2009 / Cao Qinxiang}

\begin{abstract}
Trò chơi tổ hợp là trò chơi hai người, luân phiên thao tác, thông tin hoàn
toàn. Bài viết bắt đầu từ trò chơi trừ động theo bội \(k\)
(\emph{\(k\)-based dynamic subtraction game}), sau đó bàn về những phương pháp
thường dùng để giải bài toán trò chơi tổ hợp trong tin học: định lý trạng thái
\(N/P\), quy hoạch động, đơn điệu, tìm kiếm có nhớ, hàm Sprague--Grundy, tổng
Nim, tích Nim, đối xứng và tham lam. Tác giả cũng tối ưu thuật toán cho trò
chơi trừ động từ \(\Oh(S^3)\) xuống \(\Oh(S)\), đưa ra thuật toán tính tích Nim
trong \(\Oh(\log^2 x)\), đồng thời phân tích lại lời giải chính thức của một
bài BOI 2008 để chỉ ra sai lệch trong ước lượng độ phức tạp và cải thiện về
\(\Oh(n^2)\).
\end{abstract}

\paragraph{Từ khóa.}
Trạng thái \(N/P\); đơn điệu; hàm SG; tổng Nim; tích Nim; phân tích đối xứng;
phân tích tham lam; phân tích độ phức tạp.

\section{Mở đầu}

\subsection{Định nghĩa trò chơi tổ hợp}

Từ ``trò chơi'' bao trùm rất rộng: kéo búa bao, bài, cờ tướng, cờ vây, Nim,
v.v. Trong bài viết này, một trò chơi được xem là trò chơi tổ hợp nếu thỏa
các điều kiện sau:
\begin{enumerate}
  \item Có đúng hai người chơi.
  \item Hai bên luân phiên thao tác; không có lượt thao tác đồng thời.
  \item Ở mọi thời điểm, tập thao tác hợp lệ của người đang đi là một tập hữu
  hạn và xác định; không có yếu tố ngẫu nhiên.
  \item Thông tin hoàn toàn: cả hai bên đều biết trạng thái ban đầu và toàn bộ
  lịch sử thao tác.
  \item Một số trạng thái được quy định là trạng thái kết thúc thắng, người
  nhận được chúng thắng; một số trạng thái được quy định là trạng thái kết
  thúc thua, người nhận được chúng thua. Bài viết chỉ xét các trò chơi luôn
  phân thắng thua sau hữu hạn lượt, không xét hòa.
\end{enumerate}

Vì vậy trò chơi một người như đẩy hộp, trò có thao tác đồng thời như kéo búa
bao, trò có gieo xúc xắc, hoặc trò có thông tin ẩn đều nằm ngoài phạm vi này.

Một trò chơi tổ hợp thường được mô tả bằng ánh xạ
\[
  f:X\to \mathcal{P}(X),
\]
trong đó \(X\) là tập trạng thái, \(\mathcal{P}(X)\) là tập các tập con của
\(X\), và \(f(x)\) là tập trạng thái có thể đạt được sau một nước đi từ
\(x\). Ta gọi \(y\in f(x)\) là một hậu duệ trực tiếp của \(x\). Nói cách khác,
trò chơi có thể được hiểu như một đồ thị có hướng không chu trình
\(G=(V,E)\), với \(V=X\) và
\[
  E=\{(x,y):y\in f(x)\}.
\]

\subsection{Những câu hỏi cần trả lời}

Với một trò chơi hai người, thông thường cần trả lời hai câu hỏi: ai có chiến
lược thắng, và chiến lược thắng đó là gì. Trong bài toán tin học, yêu cầu
tương ứng có ba mức:
\begin{enumerate}
  \item Chỉ xác định ai có chiến lược thắng.
  \item Tìm một nước đi thắng trong một lượt.
  \item Duy trì nước đi thắng qua nhiều lượt, như trong bài tương tác.
\end{enumerate}

\subsection{Thước đo độ phức tạp}

Với bài toán trò chơi, ngoài kích thước dữ liệu vào, còn có ba đại lượng tự
nhiên: \(S(x)\) là số hậu duệ trực tiếp hoặc gián tiếp của trạng thái hiện
tại, \(D(x)\) là lượng dữ liệu cần để mô tả trạng thái hiện tại, và \(T(x)\) là
số lượt tối đa còn lại trong trường hợp xấu nhất. Vì mô tả trạng thái phải
phân biệt được nó với các trạng thái khác, thường có
\[
  D(x)=\Theta(\log S(x)).
\]
Mặt khác, do đồ thị trò chơi không có chu trình, với mọi \(y\in f(x)\) ta có
\[
  S(x)\ge S(y)+1,
\]
nên \(S(x)=\Omega(T(x))\). Có trò chơi đạt \(S(x)=\Theta(T(x))\), chẳng hạn
trò lấy đá đơn giản. Nhưng nhiều trò chơi không như vậy: trong trò chơi trừ
động theo bội \(k\), \(S(x)=\Theta(T(x)^2)\); trong Nim \(k\) đống,
\(S(x)=\Theta(T(x)^k)\); trong Green Hackenbush,
\(S(x)=\Theta(2^{T(x)})\).

Do đó bài viết thường dùng \(T(e)\), với \(e\) là trạng thái ban đầu, làm
thước đo thời gian tự nhiên của trò chơi. Với bài tương tác, thuật toán trực
tuyến lý tưởng có dạng \(\Oh(T(e))-\Oh(1)\): tiền xử lý tuyến tính theo thời
lượng trò chơi, sau đó mỗi lượt trả lời hằng số.

\section{Bài toán xuất phát}

\paragraph{Trò chơi trừ động theo bội \(k\).}
Cho số nguyên \(S\ge 2\). Người đi trước trừ khỏi \(S\) một số \(x\), ít nhất
là \(1\) và nhỏ hơn \(S\). Sau đó hai bên luân phiên trừ khỏi \(S\) một số
nguyên dương, nhưng số bị trừ ở lượt hiện tại không được vượt quá \(k\) lần số
đối thủ vừa trừ ở lượt trước. Ai làm \(S\) trở thành \(0\) thì thắng. Hỏi ai
có chiến lược thắng.

Ta biểu diễn trạng thái bằng cặp có thứ tự \((m,n)\): hiện còn \(m\), và người
đang đi được trừ nhiều nhất \(n\). Trạng thái ban đầu là \((S,S-1)\).

Với \(k=1\) và \(k=2\), bài viết đội tuyển quốc gia năm 2002 đã cho kết quả
đóng. Khi \(k=1\), hậu thủ thắng khi và chỉ khi \(S\) là lũy thừa của \(2\).
Khi \(k=2\), hậu thủ thắng khi và chỉ khi \(S\) là một số Fibonacci. Việc kiểm
tra hai điều kiện này có thể làm trong \(\Oh(\log S)\), hoặc
\(\Oh(\log^2 S)\) nếu tính cả số học nhiều chữ số. Phần sau trình bày cách xử
lý \(k\) tổng quát và cách duy trì chiến lược thắng.

\section{Nền tảng: trạng thái N/P và quy hoạch động}

\subsection{Khái niệm trạng thái N và P}

Trạng thái \(N\) là trạng thái mà người sắp đi có chiến lược thắng
(\emph{Next player wins}). Trạng thái \(P\) là trạng thái mà người vừa đi có
chiến lược thắng (\emph{Previous player wins}).

Ví dụ, trong trò lấy đá có \(21\) viên, mỗi lượt lấy \(1\), \(2\) hoặc \(3\)
viên, ai lấy viên cuối thắng. Trạng thái còn \(0\) viên là \(P\). Các trạng
thái còn \(1,2,3\) viên là \(N\), vì người sắp đi có thể lấy hết. Trạng thái
còn \(4\) viên là \(P\), vì mọi nước đi đều đưa đối thủ tới một trạng thái
\(N\). Suy ra các trạng thái \(0,4,8,12,\ldots\) là \(P\); trạng thái ban đầu
\(21\) viên là \(N\), nên người đi trước thắng.

\paragraph{Định lý trạng thái \(N/P\).}
Một trạng thái là \(P\) nếu nó là trạng thái kết thúc thắng hoặc mọi hậu duệ
của nó đều là \(N\). Một trạng thái là \(N\) nếu nó là trạng thái kết thúc
thua hoặc có ít nhất một hậu duệ là \(P\).

\subsection{Quy hoạch động tổng quát}

Từ định lý trên, có thể tính nhãn \(N/P\) bằng quy hoạch động trên thứ tự topo
ngược của đồ thị trò chơi:
\begin{enumerate}
  \item Đánh dấu trạng thái kết thúc thắng là \(P\), trạng thái kết thúc thua
  là \(N\).
  \item Với trạng thái chưa biết mà mọi hậu duệ đã biết đều là \(N\), đánh
  dấu nó là \(P\).
  \item Với trạng thái chưa biết có một hậu duệ là \(P\), đánh dấu nó là \(N\).
  \item Nếu hai bước trước không tạo nhãn mới thì dừng; ngược lại lặp lại.
\end{enumerate}

Với đồ thị \(G(V,E)\), cách làm này có độ phức tạp \(\Oh(|E|)\). Đây là
phương pháp phổ quát nhưng thường không đủ nhanh. Với trò chơi trừ động theo
bội \(k\), ta có \(|E|=\Theta(S^3)\) và \(T(e)=\Theta(S)\), nên thuật toán
\(\Oh(S^3)\) là quá lớn. Phần tiếp theo tối ưu nó về tuyến tính theo \(S\).

\section{Tối ưu hóa dựa trên quy hoạch động}

\subsection{Đơn điệu trạng thái trong trò chơi trừ động}

Đặt
\[
\NP(m,n)=
\begin{cases}
  1, & (m,n)\text{ là trạng thái }N,\\
  0, & (m,n)\text{ là trạng thái }P.
\end{cases}
\]
Trong trò chơi trừ động, \(\NP(m,n)\) đơn điệu không giảm theo \(n\). Nếu
\(\NP(m_0,n_0)=1\), thì với mọi \(n>n_0\), trạng thái \((m_0,n)\) cũng là
\(N\), vì người chơi có ít nhất toàn bộ các nước đi hợp lệ ở \((m_0,n_0)\) và
có thể đi tới cùng trạng thái thắng.

Do đó chỉ cần lưu
\[
  f(m)=\min\{n:\NP(m,n)=1\}.
\]
Vì \(\NP(m,m)=1\), giá trị này luôn tồn tại. Theo định nghĩa, hàng \(m\) của
bảng trạng thái có dạng: các cột \(n<f(m)\) là \(P\), các cột \(n\ge f(m)\)
là \(N\).

Khi từ \((m,n)\) trừ \(r\), trạng thái mới là \((m-r,kr)\). Do đó
\((m,n)\) là \(N\) khi tồn tại \(1\le r\le n\) sao cho \((m-r,kr)\) là \(P\).
Suy ra công thức chuyển:
\[
  f(m)=\min\{n:f(m-n)>kn\},
\]
với quy ước \(f(0)=+\infty\), vì còn \(0\) là trạng thái \(P\) bất kể giới hạn
trừ ở lượt sau. Kiểm tra trực tiếp từng \(n\) cho mỗi \(m\) tạo thuật toán
\(\Oh(S^2)\). Hậu thủ thắng ở trạng thái ban đầu khi và chỉ khi \(f(S)=S\).

\subsection{Đơn điệu quyết định và ngăn xếp}

Thuật toán \(\Oh(S^2)\) vẫn chưa khai thác hết cấu trúc. Nếu vẽ bảng
\(\NP(m,n)\) và tô đen các ô \(P\), mỗi hàng tạo một đoạn liên tiếp giống như
một bức tường. Công thức
\[
  f(m)=\min\{n:f(m-n)>kn\}
\]
có thể hiểu hình học như sau: từ ô \((m,0)\), kẻ một đường thẳng có hệ số góc
\(k-1\); \(f(m)\) là khoảng cách tới bức tường đầu tiên mà đường này gặp.

Các đường thẳng này song song nhau và khi \(m\) tăng thì chúng dịch dần xuống
phải. Mỗi bức tường cố định và có đầu phải hữu hạn. Vì thế nếu một bức tường
không chặn được đường hiện tại, nó cũng sẽ không chặn được bất kỳ đường nào
sau này. Ta dùng ngăn xếp để lưu các bức tường còn hữu ích:
\begin{enumerate}
  \item Khi tính \(f(m)\), kiểm tra lần lượt bức tường trên đỉnh ngăn xếp.
  \item Nếu nó không chặn được đường xuất phát từ \((m,0)\), loại nó khỏi ngăn
  xếp.
  \item Nếu nó chặn được, dừng và suy ra \(f(m)\).
  \item Từ \(f(m)\), tạo bức tường mới nếu có, rồi đưa vào ngăn xếp.
\end{enumerate}

Mỗi bức tường được đưa vào ngăn xếp đúng một lần và bị loại nhiều nhất một
lần. Vì mỗi \(m\) tạo nhiều nhất một bức tường, tổng thời gian là \(\Oh(S)\).
Đây là thuật toán tuyến tính theo \(T(e)\).

\subsection{Tìm và duy trì chiến lược thắng}

Sau khi biết toàn bộ \(f(m)\), nếu trạng thái hiện tại \((m,n)\) là \(N\), tức
\(n\ge f(m)\), thì trừ đúng \(f(m)\) là một nước thắng, đồng thời là nước thắng
nhỏ nhất. Vì vậy trò chơi trực tuyến có thể xử lý trong
\(\Oh(S)-\Oh(1)\).

Với \(k=1\) hoặc \(k=2\) có thể có thuật toán nhanh hơn nhờ công thức số học
riêng, nhưng phần cốt lõi của trường hợp tổng quát là đơn điệu trạng thái và
đơn điệu quyết định trong quy hoạch động.

\subsection{Tìm kiếm có nhớ}

Quy hoạch động có thể triển khai bằng truy hồi có nhớ thay vì duyệt topo. Từ
định lý \(N/P\),
\[
\NP(x)=
\begin{cases}
  P, & \forall y\in f(x),\ \NP(y)=N,\\
  N, & \exists y\in f(x),\ \NP(y)=P.
\end{cases}
\]
Khi đang tìm kiếm, chỉ cần gặp một hậu duệ \(P\) là biết trạng thái hiện tại
là \(N\), không cần xét các hậu duệ còn lại. Hai chiến lược thường dùng là:
\begin{enumerate}
  \item Dùng heuristic để thử trước nước đi có nhiều khả năng là nước thắng,
  hoặc thử trước trạng thái con dễ phán đoán.
  \item Dùng ngẫu nhiên để xáo thứ tự nước đi, tránh dữ liệu đặc biệt làm tìm
  kiếm rơi vào nhánh xấu.
\end{enumerate}

Tìm kiếm có nhớ có thể tránh nhiều trạng thái không xuất hiện trong thực tế và
được lợi nhờ cắt nhánh. Ngược lại, nhiều tối ưu dựa trên đơn điệu quyết định
chỉ thuận tiện với cách tính lặp.

\subsection{Quan sát quy luật: trò knight của CEOI 2008}

Trò knight của CEOI 2008 đặt \(K\) quân mã trên bàn cờ \(N\times N\). Mỗi lượt,
người chơi phải di chuyển tất cả các quân mã còn di chuyển được theo một trong
bốn hướng hợp lệ đã cho; quân nào không đi được thì đứng yên. Ai tới lượt mà
không quân nào đi được thì thua.

Nếu chỉ có một quân mã ở \((x,y)\), vì \(x+y\) giảm sau mỗi nước, có thể dùng
quy hoạch động \(\Oh(N^2)\). Nhưng \(N\) lớn nên cần quan sát bảng nhỏ để rút
ra công thức. Quy luật của bảng \(N/P\) phụ thuộc vào \(N\bmod 4\), nhưng sau
khi phân lớp theo phần dư, có thể xác định thắng thua trong \(\Oh(1)\).

Với \(K\) quân mã, chỉ biết từng quân thắng hay thua chưa đủ. Người chơi muốn
thắng ở quân thắng càng muộn càng tốt, và thua ở quân thua càng sớm càng tốt.
Gọi \(\operatorname{Round}(x,y)\) là số lượt tới khi quân ở \((x,y)\) phân
thắng thua, và đặt
\[
RR(x,y)=
\begin{cases}
  \operatorname{Round}(x,y), & (x,y)\text{ là }N,\\
 -\operatorname{Round}(x,y), & (x,y)\text{ là }P.
\end{cases}
\]
Nếu \((x^*,y^*)\) chạy qua các hậu duệ của \((x,y)\), công thức chuyển có dạng
\[
RR(x,y)=
\begin{cases}
  -\min RR(x^*,y^*)+1, & (x,y)\text{ là }N,\\
  -\min RR(x^*,y^*)-1, & (x,y)\text{ là }P.
\end{cases}
\]
Quân có \(|RR|\) lớn nhất quyết định kết quả chung: giá trị dương nghĩa là
tiên thủ thắng, giá trị âm nghĩa là hậu thủ thắng.

Bảng thực nghiệm cho thấy
\[
  \lim_{x,y\to\infty}\frac{\operatorname{Round}(x,y)}{x+y}=\frac12,
\]
và sau khi trừ \(\lfloor x/2\rfloor+\lfloor y/2\rfloor\), phần lõi lặp theo
khối \(4\times 4\):
\[
\begin{pmatrix}
 2&-1& 1& 0\\
 1& 0& 0& 1\\
 0& 1& 0&-1\\
 1& 0& 1& 2
\end{pmatrix}.
\]
Lớp ngoài cùng vẫn phụ thuộc \(N\bmod 4\). Có thể viết công thức \(\Oh(1)\),
nhưng cách thực dụng hơn là dùng công thức cho phần trong và để quy hoạch động
xử lý lớp ngoài, đạt \(\Oh(K+N)\), đủ tốt khi \(K\) cùng bậc với \(N\).

\section{Hàm SG, tổng Nim và tích Nim}

\subsection{Hàm SG và tổng Nim}

Với tập hữu hạn \(S\) gồm các số nguyên không âm, đặt
\[
  \mex(S)=\min\{x:x\notin S,\ x\ge 0\}.
\]
Với trò chơi không có trạng thái kết thúc thua, hàm Sprague--Grundy được định
nghĩa bởi
\[
\SG(x)=
\begin{cases}
  0, & x\text{ là trạng thái kết thúc thắng},\\
  \mex\{\SG(y):y\in f(x)\}, & \text{trường hợp khác}.
\end{cases}
\]
Từ định lý \(N/P\), \(\SG(x)=0\) khi và chỉ khi \(x\) là trạng thái \(P\).

Tổng Nim là phép xor trên các số nguyên không âm, ký hiệu \(\xor\). Nó có các
tính chất giao hoán, kết hợp, \(a\xor a=0\), \(a\xor 0=a\), nên phương trình
\(a\xor x=b\) có nghiệm duy nhất \(x=a\xor b\).

\paragraph{Định lý SG.}
Nếu trò chơi tổng
\[
  G=G_1+G_2+\cdots+G_n
\]
cho phép mỗi lượt chọn đúng một trò con và đi một nước trong đó, thì với
\(x=(x_1,x_2,\ldots,x_n)\),
\[
  \SG(x)=\SG(x_1)\xor\SG(x_2)\xor\cdots\xor\SG(x_n).
\]

Định lý này biến việc tính trên tích Descartes khổng lồ thành việc tính từng
trò con. Với Nim cổ điển có \(k\) đống \(x_1,\ldots,x_k\), trò một đống có
\(\SG(x)=x\). Do đó tiên thủ thắng khi và chỉ khi
\[
  x_1\xor x_2\xor\cdots\xor x_k\ne 0.
\]

\subsection{Chiến lược thắng trong Nim và duy trì trực tuyến}

Giả sử \(X=x_1\xor\cdots\xor x_k\ne 0\). Chọn một đống \(i\) mà bit cao nhất
của \(X\) cũng bật trong \(x_i\), rồi đặt
\[
  y_i=x_i\xor X,\qquad y_j=x_j\ (j\ne i).
\]
Khi đó \(y_i<x_i\) và xor của các đống mới bằng \(0\), nên đây là nước thắng.

Nếu mỗi lượt đều tính lại từ đầu, chi phí là \(\Oh(k)\). Trong bài tương tác,
có thể duy trì tổng xor trong \(\Oh(1)\) sau mỗi lượt, vì nhiều nhất hai đống
thay đổi giữa hai lượt của ta. Để tìm nhanh đống chứa bit cần thiết, dùng
\(\log(\max x_i)\) danh sách liên kết hai chiều, mỗi danh sách lưu những đống
có bit tương ứng bằng \(1\). Khi cập nhật một đống, sửa các danh sách theo các
bit của nó, đạt \(\Oh(\log\max x_i)\) mỗi lượt sau tiền xử lý
\(\Oh(k\log\max x_i)\).

\subsection{Lồng ghép định lý SG}

Định lý SG có thể được dùng lồng nhau. Xét trò có \(n\) quân cờ xếp thành hàng;
mỗi lượt lấy một quân hoặc hai quân kề nhau, ai lấy quân cuối thắng. Nếu lấy ở
giữa, hàng cờ tách thành hai trò độc lập. Vì vậy
\[
\SG(n)=\mex\{\SG(n-1),\SG(n-2),
\SG(i)\xor\SG(n-i-1),\SG(i)\xor\SG(n-i-2)\}.
\]
Với \(n>0\), thực ra \(\SG(n)>0\): người đi trước chỉ cần lấy ở giữa để tạo hai
nửa đối xứng. Ví dụ này nhắc rằng không nên bó buộc suy nghĩ vào việc tính
toàn bộ bảng \(N/P\) hoặc SG nếu có một lập luận cấu trúc đơn giản hơn.

\subsection{Tổng Nim bậc N}

Trong Nim bậc \(N\), mỗi lượt được chọn từ \(1\) đến \(N\) đống không rỗng và
lấy tùy ý một số dương ở mỗi đống được chọn. Hậu thủ thắng khi và chỉ khi ở
mỗi bit nhị phân, số đống có bit đó bằng \(1\) là bội của \(N+1\). Kết luận
này tương tự định lý SG, nhưng khi \(N>1\) nó chỉ cho tiêu chuẩn thắng thua và
nước thắng; nó không cho một giá trị SG duy nhất để có thể lồng tiếp.

\subsection{Tích Nim từ trò lật đồng xu}

Trò lật đồng xu I đặt đồng xu trên bảng tọa độ \((0,0)\) tới \((n,m)\). Mỗi
lượt lật hai đồng xu cùng hàng hoặc cùng cột, trong đó đồng xu có tọa độ lớn
hơn phải đang ngửa và bị lật úp. Khi chỉ có một đồng xu ngửa tại \((x,y)\),
trò tương đương tổng của hai đống Nim, nên
\[
  \SG(x,y)=x\xor y.
\]

Trò lật đồng xu II thay thao tác bằng việc lật bốn đỉnh của một hình chữ nhật:
\[
  (x,y),\ (x,b),\ (a,y),\ (a,b),\qquad 0\le a<x,\ 0\le b<y,
\]
trong đó \((x,y)\) phải đang ngửa và bị lật úp. Khi chỉ có một đồng xu ngửa,
hàm SG phức tạp hơn. Ta định nghĩa tích Nim bằng
\[
  x\nimprod y
  =\mex\{(a\nimprod y)\xor(x\nimprod b)\xor(a\nimprod b):
       0\le a<x,\ 0\le b<y\}.
\]
Đây là giá trị SG của một đồng xu ở \((x,y)\) trong trò II.

\paragraph{Định lý Tartan.}
Với tích trò chơi \(G=G_1\cdot G_2\cdots G_n\) được định nghĩa sao cho một lượt
chọn đồng thời hậu duệ trong một số tọa độ và tổng hóa tất cả trạng thái sinh
ra theo quy tắc tích, ta có
\[
  \SG(x_1,\ldots,x_n)=
  \SG(x_1)\nimprod\SG(x_2)\nimprod\cdots\nimprod\SG(x_n).
\]
Phép \(\nimprod\) có đơn vị \(1\), giao hoán, kết hợp, và phân phối trên
\(\xor\):
\[
  (x\xor y)\nimprod z=(x\nimprod z)\xor(y\nimprod z).
\]

\subsection{Tính tích Nim trong thời gian bình phương log}

Tính trực tiếp từ định nghĩa mex rất chậm. Thuật toán nhanh dùng các tính chất
sau. Với \(x,y<2^{2^\alpha}\):
\[
  x\nimprod 2^{2^\alpha}=2^{2^\alpha}x,\qquad
  x\nimprod y<2^{2^\alpha},\qquad
  2^{2^\alpha}\nimprod 2^{2^\alpha}=\frac32\,2^{2^\alpha}.
\]

Ví dụ:
\[
\begin{aligned}
24\nimprod17
&=(16\xor8)\nimprod(16\xor1)\\
&=(16\nimprod16)\xor(16\nimprod8)\xor(16\nimprod1)\xor(8\nimprod1)\\
&=24\xor128\xor16\xor8=128.
\end{aligned}
\]
Một ví dụ khác:
\[
\begin{aligned}
8\nimprod8
&=(2\nimprod4)\nimprod(2\nimprod4)\\
&=(2\nimprod2)\nimprod(4\nimprod4)
=3\nimprod6\\
&=(1\xor2)\nimprod(4\xor2)
=4\xor2\xor8\xor3=13.
\end{aligned}
\]

Thuật toán đệ quy gồm hai hàm. Hàm
\(\operatorname{NimMulti}(x,y)\) tính \(x\nimprod y\), còn
\(\operatorname{NimMultiPower}(x,y)\) xử lý riêng trường hợp \(x\) là lũy thừa
của \(2\).

\paragraph{\(\operatorname{NimMulti}(x,y)\).}
\begin{enumerate}
  \item Nếu \(x<y\), đổi chỗ hai đối số.
  \item Nếu \(x<16\), tra bảng.
  \item Tìm \(a\) sao cho
  \(2^{2^a}\le x<2^{2^{a+1}}\), đặt \(M=2^{2^a}\).
  \item Viết \(x=pM+q,\ y=sM+t\), với \(0\le q,t<M\).
  \item Tính
  \[
    c_1=\operatorname{NimMulti}(p,s),\quad
    c_2=\operatorname{NimMulti}(p,t)\xor\operatorname{NimMulti}(q,s),\quad
    c_3=\operatorname{NimMulti}(q,t).
  \]
  \item Trả về
  \[
    (c_1\xor c_2)M\xor c_3
    \xor \operatorname{NimMultiPower}(M/2,c_1).
  \]
\end{enumerate}

\paragraph{\(\operatorname{NimMultiPower}(x,y)\).}
\begin{enumerate}
  \item Nếu \(x<16\), tra bảng.
  \item Tìm \(a\) sao cho \(2^{2^a}\le x<2^{2^{a+1}}\), đặt \(M=2^{2^a}\).
  \item Viết \(x=PM,\ y=SM+T\).
  \item Tính
  \[
    d_1=\operatorname{NimMultiPower}(P,S),\qquad
    d_2=\operatorname{NimMultiPower}(P,T).
  \]
  \item Trả về
  \[
    M(d_1\xor d_2)\xor\operatorname{NimMultiPower}(M/2,d_1).
  \]
\end{enumerate}

Gọi \(f(a)\) là thời gian xấu nhất của \(\operatorname{NimMulti}\) khi
\(0\le x,y<2^{2^{a+1}}\), và \(g(a)\) là thời gian của hàm phụ trong cùng miền.
Ta có
\[
  f(a)=4f(a-1)+g(a-1),\qquad g(a)=3g(a-1).
\]
Suy ra \(g(a)=\Oh(3^a)\) và \(f(a)=\Oh(4^a)\). Vì
\(a=\Theta(\log\log\max(x,y))\), thời gian là
\[
  \Oh(4^{\log\log x})=\Oh(\log^2 x).
\]

Tính đúng của hai công thức trả về đến từ phân phối của \(\nimprod\) trên
\(\xor\) và các tính chất của lũy thừa \(2^{2^a}\). Cụ thể,
\[
\begin{aligned}
\operatorname{NimMultiPower}(x,y)
&=PM\nimprod(SM+T)\\
&=(P\nimprod M)\nimprod((S\nimprod M)\xor T)\\
&=\bigl((M\nimprod M)\nimprod(P\nimprod S)\bigr)
  \xor \bigl(M\nimprod(P\nimprod T)\bigr)\\
&=M(d_1\xor d_2)\xor \operatorname{NimMultiPower}(M/2,d_1),
\end{aligned}
\]
và phép biến đổi tương tự cho \(\operatorname{NimMulti}\).

\section{Phân tích không dựa hoàn toàn vào quy hoạch động}

\subsection{Đối xứng hóa}

Đối xứng hóa là tìm cách đưa tình thế về dạng hai bên đối xứng; khi đối thủ có
nước đi, ta cũng có nước đáp lại. Trong bài US Open 2008 The Ultimate Battle,
chiến trường là lưới \(W\times H\), hai người đứng ở hai ô khác hàng và khác
cột. Mỗi lượt đi ngang hoặc dọc tùy ý nhưng không được vượt qua hàng hoặc cột
của đối thủ, vì khi cùng hàng hoặc cùng cột thì đối thủ có thể bắn.

Điều kiện thắng có thể phân tích như Nim hai đống, nhưng cách nhìn đơn giản
hơn là duy trì sau lượt của mình khoảng cách theo hàng và theo cột bằng nhau.
Nếu đối thủ lùi theo một trục, ta tiến cùng số ô trên trục đó; nếu đối thủ
tiến theo một trục, ta tiến cùng số ô trên trục kia. Vì ta luôn tiến, cuối
cùng đối thủ hết nước hợp lệ.

\subsection{Tham lam và bài BOI 2008 Game}

Trong trò BOI 2008 Game, bàn cờ \(n\times n\) có ô trắng đi được và ô đen là
chướng ngại. Hai ô \(A,B\) là vị trí xuất phát của tiên thủ và hậu thủ. Mỗi
lượt người chơi đi một ô lên, xuống, trái hoặc phải vào ô trắng. Nếu đi đúng
vào ô đối thủ đang đứng, người đó được đi thêm một bước, tức ``nhảy qua'' đối
thủ. Ai đi tới ô xuất phát của đối phương trước thì thắng, kể cả khi tới đó
nhờ bước nhảy qua.

Nếu mô tả trạng thái trực tiếp bằng \((x_1,y_1,x_2,y_2)\) và lượt đi, có
\(\Oh(n^4)\) trạng thái, quá lớn. Dấu hiệu tham lam là chữ ``trước'' trong
điều kiện thắng: cả hai bên đều muốn tới điểm xuất phát của đối thủ càng nhanh
càng tốt. Vì quãng đường ngắn nhất giữa \(A\) và \(B\) là như nhau cho cả hai,
nếu không có luật nhảy qua thì tiên thủ thắng. Vậy hậu thủ chỉ có thể thắng
nếu ép được việc nhảy qua trên một đường ngắn nhất; tiên thủ thắng nếu tránh
được điều đó.

Lời giải chính thức của BOI dùng các lớp đường ngắn nhất. Gọi \(d\) là khoảng
cách ngắn nhất giữa \(A\) và \(B\), và \(L_A[i]\) là tập ô nằm trên một đường
ngắn nhất, cách \(A\) đúng \(i\). Đặt \(\NP_A[i,j,k]\) là trạng thái tới lượt
\(A\), \(A\) đứng tại ô thứ \(j\) của \(L_A[i]\), \(B\) đứng tại ô thứ \(k\)
của \(L_A[d-i]\). Tương tự \(\NP_B[i,j,k]\) cho lượt \(B\). Lời giải chính
thức ước lượng có \(\Oh(n^3)\) trạng thái và mỗi chuyển trạng thái hằng số,
nên tổng thời gian \(\Oh(n^3)\), dùng mảng cuốn để còn \(\Oh(n^2)\) bộ nhớ.

Vấn đề là ba chỉ số hình thức không bảo đảm mỗi chiều chỉ \(\Oh(n)\). Trên một
bàn cờ có hành lang uốn lượn, \(d\) có thể là \(\Theta(n^2)\). Hơn nữa,
\(|L_A[i]|\) có thể lớn hơn tuyến tính theo những cách làm hỏng ước lượng. Tác
giả mô tả một phản ví dụ kiểu ``ống nước và bể nước'': bố trí
\(\Oh(\sqrt n)\) hàng và \(\Oh(\sqrt n)\) cột các khối, mỗi khối kích thước
\(\Oh(\sqrt n)\times\Oh(\sqrt n)\), rồi thiết kế các hành lang vào ra sao cho
rất nhiều ô trong các khối cùng thuộc một lớp khoảng cách. Khi đó có lớp
\(|L_A[i]|=\Theta(n^{1.5})\), và các lớp đối xứng từ \(B\) cũng đạt cùng bậc.
Vì vậy bộ nhớ có thể cần
\[
  \max_i |L_A[i]|\cdot |L_A[d-i]|=\Theta(n^3),
\]
còn tổng số trạng thái trên các lớp có thể đạt cận dưới \(\Omega(n^{3.5})\).
Do đó ước lượng \(\Oh(n^3)\) thời gian và \(\Oh(n^2)\) bộ nhớ của lời giải
chính thức không đúng về mặt lý thuyết.

Tuy vậy, ý tưởng tham lam vẫn đúng. Ta tô đen mọi ô trắng không nằm trên đường
ngắn nhất nào giữa \(A\) và \(B\). Với mỗi \(i\), tập \(L_A[i]\) cùng ô đen và
biên bàn cờ tạo thành một đường khép kín chia vùng trắng thành phía gần \(A\)
và phía xa \(A\). Các ô liên tiếp của \(L_A[i]\) chỉ kề với những đoạn liên
tiếp trong \(L_A[i+1]\). Vì thế, với mỗi cặp \(i,j\), tập các \(k\) làm cho
\(\NP_A[i,j,k]=N\) là một đoạn liên tiếp trên vòng \(L_A[d-i]\); với
\(\NP_B\) cũng vậy.

Thay vì lưu toàn bộ \(k\), chỉ cần lưu hai biên
\(\operatorname{left}[i,j]\) và \(\operatorname{right}[i,j]\) của đoạn \(N\)
trên vòng. Nếu \(\operatorname{left}\le\operatorname{right}\), đoạn thắng là
\([\operatorname{left},\operatorname{right}]\); nếu
\(\operatorname{left}>\operatorname{right}\), đoạn thắng vòng qua đầu mảng.
Các biên này được truy hồi đơn giản từ lớp trước. Vì mỗi ô xuất hiện trong
đúng một lớp \(L_A[i]\), tổng số cặp \((i,j)\) là \(\Oh(n^2)\), nên thời gian
và bộ nhớ đều có thể đưa về \(\Oh(n^2)\).

\section{Tổng kết}

Định lý trạng thái \(N/P\) và quy hoạch động dựa trên nó là khung phổ quát cho
bài toán trò chơi. Nhưng hiệu quả thực sự phụ thuộc vào việc phát hiện cấu
trúc riêng của bài: đơn điệu trạng thái, đơn điệu quyết định, đối xứng, tham
lam, hoặc công thức rút ra từ quan sát.

Hàm SG là phiên bản mạnh hơn của phân loại \(N/P\). Khi kết hợp với tổng Nim
và tích Nim, nó có thể xử lý nhiều trò chơi phức tạp. Tuy nhiên, khó khăn của
bài toán thường không nằm ở việc áp dụng định lý, mà ở bước nhận ra trò chơi
đã cho có thể biểu diễn thành tổng hoặc tích của các trò con như thế nào.

\appendix

\section{Trường hợp k bằng 1 của trò chơi trừ động}

Dùng ký hiệu của phần chính. Đặt \(\lowbit_2(m)\) là bit thấp nhất khác \(0\)
của \(m\) trong nhị phân:
\[
  \lowbit_2(m)=2^t,\qquad 2^t\mid m,\quad 2^{t+1}\nmid m.
\]
Khi \(k=1\), ta có
\[
  f(m)=\lowbit_2(m).
\]

Chứng minh bằng quy nạp. Với \(0<n<\lowbit_2(m)\), phép trừ nhị phân cho
\[
  \lowbit_2(m-n)=\lowbit_2(n)\le n.
\]
Theo giả thiết quy nạp,
\[
  f(m-n)=\lowbit_2(m-n)\le n=kn,
\]
nên \(f(m)\ge \lowbit_2(m)\). Mặt khác,
\[
  f(m-\lowbit_2(m))>\lowbit_2(m),
\]
vì nếu \(m=\lowbit_2(m)\) thì vế trái là \(f(0)=+\infty\); còn nếu không, xóa
bit \(1\) thấp nhất làm bit thấp nhất mới nằm ở vị trí cao hơn. Từ công thức
chuyển suy ra \(f(m)=\lowbit_2(m)\).

Vì trạng thái ban đầu thua khi và chỉ khi \(S=f(S)\), hậu thủ thắng đúng khi
\(S\) là lũy thừa của \(2\). Ở mọi trạng thái \(N\), trừ \(\lowbit_2(m)\) là
nước thắng.

\section{Trường hợp k bằng 2 của trò chơi trừ động}

Gọi \(F_0=F_1=1\), \(F_n=F_{n-1}+F_{n-2}\). Theo định lý Zeckendorf, mọi số
nguyên dương \(m\) có biểu diễn duy nhất thành tổng các số Fibonacci đôi một
khác nhau và không kề nhau:
\[
  m=\sum_{j=1}^r F_{i_j},\qquad i_j+2\le i_{j+1}.
\]
Đặt \(\lowbitF(m)\) là hạng nhỏ nhất trong biểu diễn này. Khi \(k=2\),
\[
  f(m)=\lowbitF(m).
\]

Ý tưởng chứng minh song song với trường hợp nhị phân. Nếu
\(0<n<\lowbitF(m)\), thì phần nhỏ nhất trong biểu diễn của \(m-n\) xuất phát
từ \(\lowbitF(m)-n\), và hạng Fibonacci nhỏ nhất của nó không vượt quá hai lần
\(n\):
\[
  f(m-n)\le 2n.
\]
Ngược lại, nếu \(\lowbitF(m)=F_i\), thì sau khi trừ \(F_i\), hạng nhỏ nhất còn
lại ít nhất là \(F_{i+2}>2F_i\), hoặc rơi vào \(f(0)=+\infty\). Vì vậy công
thức chuyển với \(k=2\) cho \(f(m)=\lowbitF(m)\).

Suy ra hậu thủ thắng khi và chỉ khi \(S\) là số Fibonacci. Ở trạng thái \(N\),
trừ \(\lowbitF(m)\) là nước thắng. Nếu không tính số học nhiều chữ số, việc
tìm \(\lowbitF(m)\) mất \(\Oh(\log m)\).

\section{Một vài biến thể Nim}

\begin{enumerate}
  \item Có \(n\) bậc thang, mỗi bậc có một số quân. Mỗi lượt chuyển vài quân
  từ bậc \(i\) xuống bậc \(i-1\). Bậc \(0\) là mặt đất; ai đưa mọi quân xuống
  đất thì thắng.
  \item Có các ô \(0,1,\ldots,n-1\) trên một hàng, một số quân nằm trên các ô.
  Mỗi lượt chuyển một quân sang ô có chỉ số nhỏ hơn; nhiều quân được phép cùng
  ô. Ai không đi được thì thua.
  \item Trên mỗi hàng của một bàn cờ chữ nhật có đúng một quân đen và một
  quân trắng. Người \(A\) chỉ đi quân trắng, người \(B\) chỉ đi quân đen, mỗi
  lượt đi trong cùng hàng và không được vượt qua quân khác. Ai không đi được
  thì thua.
  \item Có \(k\) đống \(x_1\le x_2\le\cdots\le x_k\). Mỗi lượt chọn một đống
  và lấy bớt đá nhưng phải giữ thứ tự không giảm, không được hoán đổi đống.
  Ai lấy viên cuối thắng.
  \item Trên một cây có gốc, mỗi lượt xóa một cạnh; phần không còn nối với gốc
  bị xóa theo. Ai xóa cạnh cuối cùng thắng.
\end{enumerate}

\section{Phác thảo chứng minh định lý SG}

Với hai trò con, cần chứng minh
\[
  \SG(x_1,x_2)=\SG(x_1)\xor\SG(x_2).
\]
Tính chất then chốt của xor là
\[
  p\xor q=\mex\{a\xor q:a<p\}\cup\{p\xor b:b<q\},
\]
hiểu theo các giá trị SG có thể đi tới từ từng trò con. Do đó mex của tập giá
trị hậu duệ trong trò tổng chính là xor của hai giá trị SG.

Với \(n\) trò, áp dụng kết quả hai trò nhiều lần:
\[
\begin{aligned}
\SG(x_1,\ldots,x_n)
&=\SG(x_1,\ldots,x_{n-1})\xor\SG(x_n)\\
&=\SG(x_1)\xor\SG(x_2)\xor\cdots\xor\SG(x_n).
\end{aligned}
\]

\section*{Tài liệu tham khảo}

\begin{enumerate}
  \item Theodore L. Turocy, Texas A\&M University, \emph{Game Theory}.
  \item Erik D. Demaine, Robert A. Hearn, \emph{Algorithmic Combinatorial Game Theory}.
  \item Thomas S. Ferguson, \emph{Game Theory}.
  \item E. R. Berlekamp, J. H. Conway, R. K. Guy, \emph{Winning Ways for Your Mathematical Plays}.
  \item Liu Rujia, Huang Liang, \emph{Nghệ thuật thuật toán và thi tin học}.
\end{enumerate}

\end{document}
